% ===========================================================================
% Gesamtskript VWL --- Musterlösungen Lektion 2
% Wettbewerb, Staatseingriffe & VGR (Rahmenplan DIHK 1.1)
% ===========================================================================

\begin{loesung}[title={Lösung Aufgabe~2.1 --- UWG oder GWB?}]
\textbf{a) Schutzzweck (4~P.: je Gesetz 2~P.)}\\
\emph{UWG:} schützt die \emph{Lauterkeit} des Wettbewerbs, also die \emph{Art und Weise},
wie Wettbewerb geführt wird. Geschützt werden Mitbewerber, Verbraucher und sonstige
Marktteilnehmer vor unlauteren geschäftlichen Handlungen (z.\,B. irreführende oder
aggressive Werbung, Herabsetzung von Mitbewerbern, unzumutbare Belästigung) sowie das
Interesse der Allgemeinheit an einem unverfälschten Wettbewerb (\S\,1 UWG).\\
\emph{GWB („Kartellgesetz``):} schützt den \emph{Bestand} bzw. die \emph{Freiheit} des
Wettbewerbs als Institution vor Beschränkungen --- durch Kartellverbot,
Missbrauchsaufsicht über marktbeherrschende Unternehmen und Fusionskontrolle. Über die
Einhaltung wacht das Bundeskartellamt.

\textbf{b) Zuordnung (6~P.: je Vorgang 2~P.)}
\begin{itemize}
  \item \textbf{(1) UWG} --- \emph{irreführende geschäftliche Handlung} (\S\,5 UWG): Die
    Angabe „Testsieger`` täuscht über die Qualität bzw. die Eigenschaften der Ware; die
    Werbung beeinflusst die Kaufentscheidung mit unwahren Angaben. Es geht um die
    \emph{Art} des Wettbewerbs, nicht um seine Beschränkung.
  \item \textbf{(2) UWG} --- \emph{unzumutbare Belästigung} (\S\,7 UWG): Telefonwerbung
    gegenüber Verbrauchern ist ohne deren vorherige ausdrückliche Einwilligung stets
    unzulässig.
  \item \textbf{(3) GWB} --- \emph{Preiskartell}, verboten nach \S\,1 GWB (europarechtlich
    Art.~101 AEUV): Die Absprache einer gemeinsamen Preiserhöhung schaltet den
    Preiswettbewerb zwischen den Anbietern aus. Dass die Vereinbarung nur mündlich auf
    einem Branchentreffen zustande käme, ist unerheblich --- auch aufeinander abgestimmte
    Verhaltensweisen sind verboten. Rheinquell sollte sich daran \emph{nicht} beteiligen
    und sich von dem Vorschlag erkennbar distanzieren.
\end{itemize}

\textbf{c) Rechtsfolgen (2~P.: je Gesetz 1~P., je zwei Nennungen)}\\
\emph{UWG:} Ansprüche auf Beseitigung und Unterlassung (\S\,8), Schadensersatz (\S\,9),
Gewinnabschöpfung (\S\,10); bei unerlaubter Telefonwerbung zusätzlich Bußgeld bis
300.000~€ (\S\,20). Klagebefugt sind u.\,a. Mitbewerber, Wirtschaftsverbände, IHK und
Verbraucherverbände.\\
\emph{GWB:} Bußgeld durch das Bundeskartellamt (bis zu 10~\% des weltweiten
Gesamtumsatzes des Vorjahres, \S\,81c GWB), Schadensersatzansprüche der Geschädigten
(\S\,33a GWB), Vorteilsabschöpfung, Unwirksamkeit der Absprache.
\end{loesung}

\begin{loesung}[title={Lösung Aufgabe~2.2 --- Kartellverbot und Ausnahmen}]
\textbf{a) Abgrenzung (2~P.)}\\
\emph{Kartell:} Die Unternehmen bleiben \emph{rechtlich selbstständig} und geben durch
die (auch mündliche) Vereinbarung nur einen \emph{Teil ihrer wirtschaftlichen}
Selbstständigkeit auf (z.\,B. bei der Preis- oder Gebietspolitik).
\emph{Fusion:} Die Unternehmen verschmelzen zu einer Einheit und verlieren ihre
\emph{rechtliche und wirtschaftliche} Selbstständigkeit vollständig. Für Fusionen gilt
nicht das Kartellverbot, sondern die Fusionskontrolle (\S\,35 GWB).

\textbf{b) Beurteilung der Vorhaben (8~P.: je Vorhaben 2~P. --- 1~P. Einordnung, 1~P. Begründung)}
\begin{itemize}
  \item \textbf{(1) Submissionskartell --- verboten.} Bei der Ausschreibung wird
    abgesprochen, welches Unternehmen zum Zug kommt; der öffentliche Auftraggeber zahlt
    überhöhte Preise. Verstoß gegen \S\,1 GWB; zusätzlich \emph{Straftat} nach \S\,298
    StGB (wettbewerbsbeschränkende Absprachen bei Ausschreibungen).
  \item \textbf{(2) Mittelstandskartell --- zulässig} (\S\,3 GWB). Kleine und mittlere
    Unternehmen rationalisieren durch zwischenbetriebliche Zusammenarbeit (gemeinsamer
    Einkauf), um Größennachteile gegenüber Großkonzernen auszugleichen. Voraussetzung:
    Der Wettbewerb wird nicht wesentlich beeinträchtigt und die Wettbewerbsfähigkeit der
    KMU verbessert --- hier erfüllt, da Preise und Angebote weiter selbstständig gestaltet
    werden.
  \item \textbf{(3) Forschungs- und Entwicklungskooperation --- freigestellt und damit
    zulässig.} Die Zusammenarbeit fördert den technischen Fortschritt, ohne den
    Wettbewerb zwischen den Beteiligten auszuschalten (getrennte Preissetzung, geringer
    gemeinsamer Marktanteil). Sie fällt unter die \emph{EU-Gruppenfreistellungsverordnung
    für F\&E-Vereinbarungen} bzw. die Legalausnahme nach \S\,2 GWB / Art.~101 Abs.~3 AEUV.
    Die Unternehmen müssen die Zulässigkeit selbst prüfen (Selbstprüfungssystem).
  \item \textbf{(4) Gebietskartell --- verboten.} Die Absatzmärkte werden regional
    aufgeteilt; in jedem Gebiet entfällt ein Wettbewerber, die Auftraggeber verlieren
    Auswahlmöglichkeiten. Verstoß gegen \S\,1 GWB.
\end{itemize}
\end{loesung}

\begin{loesung}[title={Lösung Aufgabe~2.3 --- Fusionskontrolle und Missbrauchsaufsicht}]
\textbf{a) Anmeldepflicht nach \S\,35 Abs.~1 GWB (5~P.)}\\
Anmeldepflichtig ist ein Zusammenschluss, wenn im letzten Geschäftsjahr \emph{beide}
Schwellen überschritten sind (1~P.):
\begin{enumerate}[label=(\arabic*)]
  \item weltweiter Umsatz aller beteiligten Unternehmen zusammen \emph{mehr als
    500~Mio.~€}: $430 + 85 = 515$~Mio.~€ $> 500$~Mio.~€ (erfüllt) (1~P.)
  \item im Inland mindestens ein Unternehmen \emph{mehr als 50~Mio.~€} (Nordlicht:
    390~Mio.~€, erfüllt) \emph{und} ein anderes \emph{mehr als 17,5~Mio.~€}
    (Hansemilch: 22~Mio.~€, erfüllt) (1~P.)
\end{enumerate}
$\Rightarrow$ Der Zusammenschluss ist \textbf{anmeldepflichtig} (1~P.).\\
\emph{Variante:} Bei nur 15~Mio.~€ Inlandsumsatz liegt Hansemilch unter 17,5~Mio.~€;
die zweite Inlandsschwelle ist nicht erfüllt $\Rightarrow$ \textbf{keine Anmeldepflicht}
nach den Umsatzschwellen (1~P.). (Ergänzend: Eine Anmeldepflicht könnte dann nur noch
über die Transaktionswertschwelle nach \S\,35 Abs.~1a GWB --- Gegenleistung über
400~Mio.~€ und erhebliche Inlandstätigkeit --- entstehen; das ist hier ausgeschlossen.)

\textbf{b) Verfahren und Untersagungsvoraussetzung (3~P.)}\\
Bis zur Freigabe gilt ein \emph{Vollzugsverbot} (\S\,41 GWB). Das Bundeskartellamt prüft
zunächst innerhalb eines Monats (Vorprüfverfahren); bei Bedenken leitet es das
Hauptprüfverfahren ein (Entscheidung grundsätzlich innerhalb von fünf Monaten ab
Anmeldung) (1~P.). Untersagt wird ein Zusammenschluss, wenn er \emph{wirksamen Wettbewerb
erheblich behindern} würde, insbesondere weil eine \emph{marktbeherrschende Stellung}
entsteht oder verstärkt wird (\S\,36 Abs.~1 GWB) (1~P.). Ein Unternehmen mit mindestens
40~\% Marktanteil gilt nach \S\,18 Abs.~4 GWB als marktbeherrschend (Vermutung) ---
mit 43~\% wäre diese Schwelle überschritten, eine Untersagung ist daher wahrscheinlich
(1~P.).

\textbf{c) Ministererlaubnis (2~P.)}\\
Nordlicht kann beim Bundeswirtschaftsminister eine \emph{Ministererlaubnis} (\S\,42 GWB)
beantragen. Er kann den untersagten Zusammenschluss erlauben, wenn die
gesamtwirtschaftlichen Vorteile die Wettbewerbsbeschränkung aufwiegen oder er durch ein
überragendes Interesse der Allgemeinheit gerechtfertigt ist und die marktwirtschaftliche
Ordnung nicht gefährdet wird. Die Erlaubnis ist eine seltene Ausnahme.

\textbf{d) Missbrauchsform (2~P.)}\\
\emph{Ausbeutungsmissbrauch} (\S\S\,19\,ff. GWB): Nordlicht nutzt ihre (Nachfrage-)Macht
gegenüber wirtschaftlich abhängigen Lieferanten aus, um unangemessen niedrige
Einkaufspreise durchzusetzen. Die Landwirte haben kaum Ausweichmöglichkeiten. Ein
\emph{Behinderungsmissbrauch} läge dagegen vor, wenn Konkurrenten gezielt verdrängt oder
vom Markt ferngehalten würden.
\end{loesung}

\begin{loesung}[title={Lösung Aufgabe~2.4 --- Höchstpreis auf dem Pelletmarkt}]
\textbf{a) Gleichgewicht (3~P.)}\\
Gleichsetzen: $900 - 1{,}5\,p = -300 + 2{,}5\,p \;\Rightarrow\; 1.200 = 4\,p
\;\Rightarrow\; p_0 = \textbf{300~€/t}$ (2~P.)\\
Einsetzen: $x_0 = 900 - 1{,}5 \cdot 300 = 900 - 450 = \textbf{450} \text{ (Tsd.~t)}$;
Probe: $x_A = -300 + 2{,}5 \cdot 300 = 450$ (stimmt) (1~P.)

\textbf{b) Höchstpreis 240~€/t (4~P.)}\\
$x_N = 900 - 1{,}5 \cdot 240 = 900 - 360 = \textbf{540 Tsd.~t}$ (1~P.)\\
$x_A = -300 + 2{,}5 \cdot 240 = -300 + 600 = \textbf{300 Tsd.~t}$ (1~P.)\\
\emph{Nachfrageüberhang} (Angebotslücke) von $540 - 300 = \textbf{240 Tsd.~t}$
($x_A < x_0 < x_N$) (1~P.).\\
Ein Höchstpreis von 320~€/t liegt \emph{über} dem Gleichgewichtspreis von 300~€/t. Der
Markt würde sich weiterhin beim zulässigen Gleichgewichtspreis von 300~€/t räumen --- der
Höchstpreis wäre \emph{wirkungslos} (1~P.).

\textbf{c) Schwarzmarktpreis und flankierende Maßnahmen (3~P.)}\\
Angeboten werden nur 300~Tsd.~t. Zahlungsbereitschaft für diese Menge aus der
Nachfragefunktion: $300 = 900 - 1{,}5\,p \;\Rightarrow\; 1{,}5\,p = 600
\;\Rightarrow\; p = \textbf{400~€/t}$. Auf dem Schwarzmarkt könnten also bis zu 400~€/t
erzielt werden --- mehr als der Gleichgewichtspreis (1~P.).\\
Flankierende Maßnahmen (2~P., je 1~P.):
\emph{Rationierung/Kontingentierung} der knappen Menge (z.\,B. Bezugsscheine je
Haushalt), damit die Verteilung nicht nach Zufall oder Beziehungen erfolgt;
\emph{Subventionen an die Anbieter}, damit trotz niedrigen Preises genügend produziert
wird; ggf. \emph{Kontrollen und Sanktionen} gegen den Schwarzhandel bzw. Förderung von
Importen.

\textbf{d) Umsatzveränderung der Anbieter (2~P.)}\\
Vorher: $300~\text{€/t} \cdot 450.000~\text{t} = 135~\text{Mio.~€}$;
nachher: $240~\text{€/t} \cdot 300.000~\text{t} = 72~\text{Mio.~€}$ (1~P.)\\
Veränderung: $72 - 135 = \textbf{$-$63~Mio.~€}$; prozentual
$-63 \div 135 \cdot 100 = \textbf{$-$46,67~\%}$ (1~P.)
\end{loesung}

\begin{loesung}[title={Lösung Aufgabe~2.5 --- Marktkonforme und marktkonträre Eingriffe}]
\textbf{a) Zuordnung (4~P.: je Vorschlag 1~P.)}
\begin{itemize}
  \item \textbf{(1) marktkonträr} --- staatlich festgesetzter \emph{Höchstpreis}; der
    Preis wird unmittelbar vorgegeben, der Preismechanismus außer Kraft gesetzt.
  \item \textbf{(2) marktkonform} --- Steuersenkung; sie verändert die Kosten der
    Anbieter, der Preis bildet sich weiterhin über Angebot und Nachfrage.
  \item \textbf{(3) marktkonform} --- \emph{Transferleistung} an Nachfrager; sie stärkt
    deren Kaufkraft, ohne den Preis festzulegen.
  \item \textbf{(4) marktkonform} --- \emph{Subvention} an Anbieter; sie beeinflusst das
    Angebot, der Preismechanismus bleibt erhalten.
\end{itemize}

\textbf{b) Wirkung der Stromsteuersenkung (3~P.)}\\
Die Stromsteuer ist ein Kostenbestandteil der Versorger. Sinkt sie, können die Anbieter
zu jedem Preis mehr bzw. jede Menge billiger anbieten: Die \emph{Angebotskurve verschiebt
sich nach rechts (unten)} (1~P.). Bei unveränderter Nachfrage \emph{sinkt der
Gleichgewichtspreis} (1~P.) und die \emph{Gleichgewichtsmenge steigt} (1~P.). Der Staat
nimmt weniger Stromsteuer ein; ein Teil der Entlastung kann bei den Anbietern verbleiben,
wenn sie die Senkung nicht voll weitergeben.

\textbf{c) Folgen der Preisobergrenze (3~P.: je Folge 1~P.)}
\begin{itemize}
  \item \emph{Nachfrageüberhang:} Zum niedrigen Preis wird mehr Strom nachgefragt, die
    Versorger bieten weniger an.
  \item \emph{Verluste und Rückzug von Anbietern:} Liegen die Beschaffungskosten der
    Stadtwerke über 30~Cent/kWh, machen sie mit jeder kWh Verlust; Anbieter nehmen keine
    Neukunden mehr auf oder ziehen sich aus dem Markt zurück, Investitionen (z.\,B. in
    Netze und erneuerbare Erzeugung) unterbleiben.
  \item \emph{Ausfall der Preisfunktionen:} Der Preis signalisiert die Knappheit nicht
    mehr (Signalfunktion), Sparanreize entfallen; der Staat müsste rationieren oder doch
    subventionieren (Folgekosten).
\end{itemize}

\textbf{d) Beurteilung (2~P.)}\\
Vorschlag (3) ist vorzuziehen: Der Zuschuss ist \emph{marktkonform} und
\emph{zielgenau} --- er entlastet nur bedürftige Haushalte, während der Strompreis
weiterhin die Knappheit anzeigt und Sparanreize sowie Investitionsanreize erhalten
bleiben. Der Höchstpreis (1) entlastet dagegen alle Haushalte unabhängig vom Einkommen,
verursacht Versorgungsengpässe und Folgeeingriffe. Nachteil von (3): Finanzierung durch
die Steuerzahler und Verwaltungsaufwand.
\end{loesung}

\begin{loesung}[title={Lösung Aufgabe~2.6 --- VGR: Vom BIP zum Volkseinkommen}]
\textbf{a) Außenbeitrag und BIP (4~P.)}\\
Außenbeitrag $=$ Exporte $-$ Importe $= 1.950 - 1.730 = \textbf{220~Mrd.~€}$ (1~P.)\\
BIP $=$ Private Konsumausgaben $+$ Konsumausgaben des Staates $+$ Bruttoinvestitionen
$+$ Außenbeitrag (1~P.)\\
BIP $= 2.180 + 860 + 740 + 220 = \textbf{4.000~Mrd.~€}$ (2~P.)

\textbf{b) BNE (3~P.)}\\
BNE $=$ BIP $-$ Primäreinkommen an die übrige Welt $+$ Primäreinkommen aus der übrigen Welt\\
BNE $= 4.000 - 145 + 210 = \textbf{4.065~Mrd.~€}$ (2~P.)\\
Das BNE ist um 65~Mrd.~€ größer, weil Inländer aus dem Ausland mehr Einkommen beziehen
(210~Mrd.~€) als Nichtinländer im Inland erzielen (145~Mrd.~€) --- der Saldo der
Primäreinkommen aus der übrigen Welt ist positiv (Inländerkonzept) (1~P.).

\textbf{c) Nettonationaleinkommen und Volkseinkommen (4~P.)}\\
Nettonationaleinkommen $=$ BNE $-$ Abschreibungen $= 4.065 - 690 =
\textbf{3.375~Mrd.~€}$ (2~P.)\\
Volkseinkommen $=$ Nettonationaleinkommen $-$ (Produktions- und Importabgaben $-$
Subventionen)\\
Volkseinkommen $= 3.375 - (420 - 70) = 3.375 - 350 = \textbf{3.025~Mrd.~€}$ (2~P.)

\textbf{d) Unternehmens- und Vermögenseinkommen, Lohn- und Gewinnquote (3~P.)}\\
Unternehmens- und Vermögenseinkommen $=$ Volkseinkommen $-$ Arbeitnehmerentgelt
$= 3.025 - 1.936 = \textbf{1.089~Mrd.~€}$ (1~P.)\\
Lohnquote $= 1.936 \div 3.025 \cdot 100 = \textbf{64~\%}$ (1~P.)\\
Gewinnquote $= 1.089 \div 3.025 \cdot 100 = \textbf{36~\%}$
(Kontrolle: $64~\% + 36~\% = 100~\%$) (1~P.)

\textbf{e) Steigende Lohnquote bei konstanten Löhnen (2~P.)}\\
Die Lohnquote ist ein \emph{Anteil}: Steigt sie bei unveränderten Arbeitnehmerentgelten,
muss das Volkseinkommen gesunken sein, d.\,h. die \emph{Unternehmens- und
Vermögenseinkommen sind zurückgegangen} --- typisch für einen konjunkturellen Abschwung,
in dem Gewinne stärker schwanken als Löhne (1~P.). Rechnerisch: $1.936 \div 0{,}66
\approx 2.933$~Mrd.~€ Volkseinkommen, also rund 92~Mrd.~€ weniger Unternehmens- und
Vermögenseinkommen. Eine steigende Lohnquote bedeutet also nicht automatisch, dass es
den Arbeitnehmern besser geht (1~P.).
\end{loesung}

\begin{loesung}[title={Lösung Aufgabe~2.7 --- Nominales und reales Wachstum}]
\textbf{a) Reales BIP (3~P.)}\\
Formel: reales BIP $=$ nominales BIP $\div$ Deflator $\cdot$ 100 (1~P.)\\
Jahr~2: $4.202{,}4 \div 103 \cdot 100 = \textbf{4.080,0~Mrd.~€}$ (1~P.)\\
Jahr~3: $4.263{,}0 \div 105 \cdot 100 = \textbf{4.060,0~Mrd.~€}$ (1~P.)\\
(Jahr~1 als Basisjahr: real $=$ nominal $= 4.000{,}0$~Mrd.~€)

\textbf{b) Wachstumsraten (4~P.: je Wert 1~P.)}\\
Formel: Wachstumsrate $=$ (Wert Berichtsjahr $-$ Wert Vorjahr) $\div$ Wert Vorjahr $\cdot$ 100
\begin{itemize}
  \item Jahr~2 nominal: $(4.202{,}4 - 4.000{,}0) \div 4.000{,}0 \cdot 100 = \textbf{+5,06~\%}$
  \item Jahr~2 real: $(4.080{,}0 - 4.000{,}0) \div 4.000{,}0 \cdot 100 = \textbf{+2,00~\%}$
  \item Jahr~3 nominal: $(4.263{,}0 - 4.202{,}4) \div 4.202{,}4 \cdot 100 = \textbf{+1,44~\%}$
  \item Jahr~3 real: $(4.060{,}0 - 4.080{,}0) \div 4.080{,}0 \cdot 100 = \textbf{$-$0,49~\%}$
\end{itemize}

\textbf{c) Preissteigerung und Erläuterung (3~P.)}\\
Preissteigerung Jahr~2 $\rightarrow$ Jahr~3: $(105 - 103) \div 103 \cdot 100 =
\textbf{1,94~\%}$ (1~P.). (Kontrolle: $1{,}0144 \approx 0{,}9951 \cdot 1{,}0194$.)\\
Erläuterung (2~P.): Das \emph{nominale} BIP steigt in Jahr~3 nur deshalb um 1,44~\%, weil
die Preise um 1,94~\% gestiegen sind. Die produzierte Gütermenge --- gemessen am
\emph{realen} BIP --- ist dagegen um 0,49~\% \emph{gesunken}. Wirtschaftswachstum wird
stets am realen BIP gemessen; die Medien sprechen daher zu Recht von einer schrumpfenden
Wirtschaft. Das Wachstumsziel des magischen Vierecks (rund 2~\% reales Wachstum) ist in
Jahr~3 verfehlt, in Jahr~2 dagegen erreicht.

\textbf{d) BIP-Deflator vs. VPI (2~P.)}\\
Der \emph{BIP-Deflator} misst die Preisentwicklung \emph{aller im Inland produzierten}
Güter und Dienstleistungen (Konsum, Investitionen, Staatsleistungen, Exporte);
Importgüter sind nicht enthalten (1~P.). Der \emph{Verbraucherpreisindex} misst die
Preisentwicklung eines repräsentativen \emph{Warenkorbs der privaten Haushalte} ---
einschließlich importierter Konsumgüter (z.\,B. Energie), aber ohne Investitionsgüter.
Er dient zur Messung der Inflationsrate (Ziel Preisniveaustabilität) (1~P.).
\end{loesung}

\begin{loesung}[title={Lösung Aufgabe~2.8 --- Magisches Viereck: Lage beurteilen}]
\textbf{a) Zielerreichung (8~P.: je Ziel 2~P. --- Zielwert und Beurteilung)}
\begin{itemize}
  \item \emph{Preisniveaustabilität} --- Zielwert: Inflationsrate von 2~\% (HVPI,
    symmetrisches EZB-Ziel). Mit 4,1~\% deutlich überschritten $\Rightarrow$
    \textbf{nicht erreicht}; die Kaufkraft sinkt spürbar.
  \item \emph{Hoher Beschäftigungsstand} --- Zielwert: Arbeitslosenquote von ca.
    3--4~\%. Mit 3,4~\% \textbf{erreicht}; es besteht nur noch Sockelarbeitslosigkeit.
  \item \emph{Stetiges und angemessenes Wirtschaftswachstum} --- Zielwert: reales
    BIP-Wachstum von rund 2~\% p.\,a. Mit 0,2~\% nahezu Stagnation $\Rightarrow$
    \textbf{nicht erreicht}.
  \item \emph{Außenwirtschaftliches Gleichgewicht} --- Zielwert: positiver Außenbeitrag
    von 1--2~\% des nominalen BIP. Mit 5,5~\% deutlich darüber $\Rightarrow$
    \textbf{nicht erreicht}; auch ein zu hoher Überschuss ist ein Ungleichgewicht.
\end{itemize}
Fazit: Nur ein Ziel ist erreicht --- die Lage zeigt typische Züge einer
\emph{Stagflation} (hohe Inflation bei stagnierendem Wachstum).

\textbf{b) Zielkonflikt (3~P.)}\\
Kürzungen der Staatsausgaben senken die gesamtwirtschaftliche Nachfrage (1~P.). Dies
dämpft den Preisauftrieb und dient der \emph{Preisniveaustabilität} (1~P.). Zugleich
erhalten Unternehmen weniger Aufträge, investieren und produzieren weniger: Das ohnehin
schwache \emph{Wachstum} sinkt weiter, und der \emph{hohe Beschäftigungsstand} ist
gefährdet, weil Arbeitskräfte abgebaut werden. Es besteht also ein Zielkonflikt zwischen
Preisniveaustabilität einerseits und Wachstum/Beschäftigung andererseits (1~P.).

\textbf{c) Beurteilung der Aussage (3~P.)}\\
Teilweise richtig: Ein positiver Außenbeitrag zeigt eine hohe internationale
Wettbewerbsfähigkeit der Exportwirtschaft (1~P.). Dauerhaft hohe Überschüsse sind aber
ein \emph{außenwirtschaftliches Ungleichgewicht}: Das Land produziert mehr, als es
selbst verbraucht und investiert (schwache Binnennachfrage); es entsteht eine hohe
\emph{Abhängigkeit vom Ausland} (Konjunktur- und Handelskonflikte schlagen direkt durch)
(1~P.). Den Überschüssen stehen Defizite und Verschuldung der Handelspartner gegenüber;
das kann zu Gegenmaßnahmen wie Zöllen sowie zu Aufwertungsdruck führen. Die Aussage „je
höher, desto besser`` ist daher abzulehnen --- Ziel ist ein moderater Überschuss
(1~P.).
\end{loesung}
